% Options for packages loaded elsewhere
\PassOptionsToPackage{unicode}{hyperref}
\PassOptionsToPackage{hyphens}{url}
\documentclass[
]{article}
\usepackage{xcolor}
\usepackage[margin=1in]{geometry}
\usepackage{amsmath,amssymb}
\setcounter{secnumdepth}{5}
\usepackage{iftex}
\ifPDFTeX
  \usepackage[T1]{fontenc}
  \usepackage[utf8]{inputenc}
  \usepackage{textcomp} % provide euro and other symbols
\else % if luatex or xetex
  \usepackage{unicode-math} % this also loads fontspec
  \defaultfontfeatures{Scale=MatchLowercase}
  \defaultfontfeatures[\rmfamily]{Ligatures=TeX,Scale=1}
\fi
\usepackage{lmodern}
\ifPDFTeX\else
  % xetex/luatex font selection
\fi
% Use upquote if available, for straight quotes in verbatim environments
\IfFileExists{upquote.sty}{\usepackage{upquote}}{}
\IfFileExists{microtype.sty}{% use microtype if available
  \usepackage[]{microtype}
  \UseMicrotypeSet[protrusion]{basicmath} % disable protrusion for tt fonts
}{}
\makeatletter
\@ifundefined{KOMAClassName}{% if non-KOMA class
  \IfFileExists{parskip.sty}{%
    \usepackage{parskip}
  }{% else
    \setlength{\parindent}{0pt}
    \setlength{\parskip}{6pt plus 2pt minus 1pt}}
}{% if KOMA class
  \KOMAoptions{parskip=half}}
\makeatother
\usepackage{color}
\usepackage{fancyvrb}
\newcommand{\VerbBar}{|}
\newcommand{\VERB}{\Verb[commandchars=\\\{\}]}
\DefineVerbatimEnvironment{Highlighting}{Verbatim}{commandchars=\\\{\}}
% Add ',fontsize=\small' for more characters per line
\usepackage{framed}
\definecolor{shadecolor}{RGB}{248,248,248}
\newenvironment{Shaded}{\begin{snugshade}}{\end{snugshade}}
\newcommand{\AlertTok}[1]{\textcolor[rgb]{0.94,0.16,0.16}{#1}}
\newcommand{\AnnotationTok}[1]{\textcolor[rgb]{0.56,0.35,0.01}{\textbf{\textit{#1}}}}
\newcommand{\AttributeTok}[1]{\textcolor[rgb]{0.13,0.29,0.53}{#1}}
\newcommand{\BaseNTok}[1]{\textcolor[rgb]{0.00,0.00,0.81}{#1}}
\newcommand{\BuiltInTok}[1]{#1}
\newcommand{\CharTok}[1]{\textcolor[rgb]{0.31,0.60,0.02}{#1}}
\newcommand{\CommentTok}[1]{\textcolor[rgb]{0.56,0.35,0.01}{\textit{#1}}}
\newcommand{\CommentVarTok}[1]{\textcolor[rgb]{0.56,0.35,0.01}{\textbf{\textit{#1}}}}
\newcommand{\ConstantTok}[1]{\textcolor[rgb]{0.56,0.35,0.01}{#1}}
\newcommand{\ControlFlowTok}[1]{\textcolor[rgb]{0.13,0.29,0.53}{\textbf{#1}}}
\newcommand{\DataTypeTok}[1]{\textcolor[rgb]{0.13,0.29,0.53}{#1}}
\newcommand{\DecValTok}[1]{\textcolor[rgb]{0.00,0.00,0.81}{#1}}
\newcommand{\DocumentationTok}[1]{\textcolor[rgb]{0.56,0.35,0.01}{\textbf{\textit{#1}}}}
\newcommand{\ErrorTok}[1]{\textcolor[rgb]{0.64,0.00,0.00}{\textbf{#1}}}
\newcommand{\ExtensionTok}[1]{#1}
\newcommand{\FloatTok}[1]{\textcolor[rgb]{0.00,0.00,0.81}{#1}}
\newcommand{\FunctionTok}[1]{\textcolor[rgb]{0.13,0.29,0.53}{\textbf{#1}}}
\newcommand{\ImportTok}[1]{#1}
\newcommand{\InformationTok}[1]{\textcolor[rgb]{0.56,0.35,0.01}{\textbf{\textit{#1}}}}
\newcommand{\KeywordTok}[1]{\textcolor[rgb]{0.13,0.29,0.53}{\textbf{#1}}}
\newcommand{\NormalTok}[1]{#1}
\newcommand{\OperatorTok}[1]{\textcolor[rgb]{0.81,0.36,0.00}{\textbf{#1}}}
\newcommand{\OtherTok}[1]{\textcolor[rgb]{0.56,0.35,0.01}{#1}}
\newcommand{\PreprocessorTok}[1]{\textcolor[rgb]{0.56,0.35,0.01}{\textit{#1}}}
\newcommand{\RegionMarkerTok}[1]{#1}
\newcommand{\SpecialCharTok}[1]{\textcolor[rgb]{0.81,0.36,0.00}{\textbf{#1}}}
\newcommand{\SpecialStringTok}[1]{\textcolor[rgb]{0.31,0.60,0.02}{#1}}
\newcommand{\StringTok}[1]{\textcolor[rgb]{0.31,0.60,0.02}{#1}}
\newcommand{\VariableTok}[1]{\textcolor[rgb]{0.00,0.00,0.00}{#1}}
\newcommand{\VerbatimStringTok}[1]{\textcolor[rgb]{0.31,0.60,0.02}{#1}}
\newcommand{\WarningTok}[1]{\textcolor[rgb]{0.56,0.35,0.01}{\textbf{\textit{#1}}}}
\usepackage{longtable,booktabs,array}
\newcounter{none} % for unnumbered tables
\usepackage{calc} % for calculating minipage widths
% Correct order of tables after \paragraph or \subparagraph
\usepackage{etoolbox}
\makeatletter
\patchcmd\longtable{\par}{\if@noskipsec\mbox{}\fi\par}{}{}
\makeatother
% Allow footnotes in longtable head/foot
\IfFileExists{footnotehyper.sty}{\usepackage{footnotehyper}}{\usepackage{footnote}}
\makesavenoteenv{longtable}
\usepackage{graphicx}
\makeatletter
\newsavebox\pandoc@box
\newcommand*\pandocbounded[1]{% scales image to fit in text height/width
  \sbox\pandoc@box{#1}%
  \Gscale@div\@tempa{\textheight}{\dimexpr\ht\pandoc@box+\dp\pandoc@box\relax}%
  \Gscale@div\@tempb{\linewidth}{\wd\pandoc@box}%
  \ifdim\@tempb\p@<\@tempa\p@\let\@tempa\@tempb\fi% select the smaller of both
  \ifdim\@tempa\p@<\p@\scalebox{\@tempa}{\usebox\pandoc@box}%
  \else\usebox{\pandoc@box}%
  \fi%
}
% Set default figure placement to htbp
\def\fps@figure{htbp}
\makeatother
\setlength{\emergencystretch}{3em} % prevent overfull lines
\providecommand{\tightlist}{%
  \setlength{\itemsep}{0pt}\setlength{\parskip}{0pt}}
\usepackage{float}
\restylefloat{figure}
\restylefloat{table}
\usepackage{booktabs}
\usepackage{longtable}
\usepackage{array}
\usepackage{multirow}
\usepackage{wrapfig}
\usepackage{float}
\usepackage{colortbl}
\usepackage{pdflscape}
\usepackage{tabularx}
\usepackage{xltabular}
\usepackage{threeparttable}
\usepackage{threeparttablex}
\usepackage[normalem]{ulem}
\usepackage{makecell}
\usepackage{xcolor}
\usepackage{bookmark}
\IfFileExists{xurl.sty}{\usepackage{xurl}}{} % add URL line breaks if available
\urlstyle{same}
\hypersetup{
  pdftitle={Estudo de Caso 01: Comparação do IMC médio de alunos do PPGEE-UFMG ao longo de dois semestres},
  pdfauthor={Alexandra Milaya Soto Aponte,},
  hidelinks,
  pdfcreator={LaTeX via pandoc}}

\title{Estudo de Caso 01: Comparação do IMC médio de alunos do PPGEE-UFMG ao longo de dois semestres}
\usepackage{etoolbox}
\makeatletter
\providecommand{\subtitle}[1]{% add subtitle to \maketitle
  \apptocmd{\@title}{\par {\large #1 \par}}{}{}
}
\makeatother
\subtitle{EEE933 - Planejamento e Análise de Experimentos}
\author{Alexandra Milaya Soto Aponte,}
\date{11 de setembro de 2026}

\begin{document}
\maketitle

\section{Resumo}\label{resumo}

Este relatório compara o Índice de Massa Corporal (IMC) médio de alunos de pós-graduação
do PPGEE-UFMG entre os semestres 2016-2 e 2017-2. A comparação foi conduzida em três
recortes independentes: a população feminina, a população masculina e a população agregada
(masculina + feminina). Em cada caso foram verificadas as premissas do teste t de Student
para duas amostras independentes (normalidade, homogeneidade de variâncias e independência
das observações) e, em seguida, aplicado o teste bilateral com nível de confiança de 95\%.
Nos três casos o teste falhou em rejeitar a hipótese nula (p-valores de 0,2612; 0,7732 e
0,7370, respectivamente).

\section{Projeto experimental}\label{projeto-experimental}

\subsection{Objetivo do Experimento e Hipóteses a serem Testadas}\label{objetivo-do-experimento-e-hipuxf3teses-a-serem-testadas}

O objetivo principal deste experimento é verificar se existe diferença estatisticamente significativa no IMC médio dos alunos de pós-graduação do PPGEE-UFMG entre os semestres 2016-2 e 2017-2. Para tanto, as amostras foram separadas em subcategorias de análise classificadas por sexo e semestre:

\begin{itemize}
\tightlist
\item
  Homens em 2016-2 vs.~Homens em 2017-2
\item
  Mulheres em 2016-2 vs.~Mulheres em 2017-2
\end{itemize}

Para cada subpopulação (separada por sexo), definem-se as seguintes hipóteses bilaterais, por não haver razão teórica prévia para supor a direção de uma eventual diferença do IMC médio da população de 2016-2 em relação à população de 2017-2:

\begin{itemize}
\tightlist
\item
  \textbf{Hipótese Nula (\(H_0\)):} Não há diferença na média do IMC entre os alunos de 2016-2 e 2017-2.
\item
  \textbf{Hipótese Alternativa (\(H_1\)):} Há diferença significativa na média do IMC entre os semestres 2016-2 e 2017-2.
\end{itemize}

\[
\begin{cases}
H_0: \mu_{2016} - \mu_{2017} = 0 \\
H_1: \mu_{2016} - \mu_{2017} \neq 0
\end{cases}
\]

\subsection{Tipos de Teste Utilizados}\label{tipos-de-teste-utilizados}

Para a comparação simples entre duas populações independentes, é necessário verificar previamente duas premissas essenciais:

\begin{enumerate}
\def\labelenumi{\arabic{enumi}.}
\tightlist
\item
  Normalidade;
\item
  Igualdade de variâncias.
\end{enumerate}

A premissa de normalidade será avaliada por meio do teste de Shapiro-Wilk. Enquanto a igualdade de variâncias será verificada utilizando o teste de Fligner-Killeen. A validação dessas premissas determinará a escolha do teste de hipótese para a comparação das médias do IMC entre os semestres:

\begin{enumerate}
\def\labelenumi{\arabic{enumi}.}
\tightlist
\item
  Teste \(t\) de Student tradicional: aplicado caso as variâncias sejam homogêneas.
\item
  Teste \(t\) de Welch: aplicado caso sejam identificadas variâncias distintas.
\end{enumerate}

É comum, na prática, optar-se diretamente pelo teste de Welch, uma vez que essa escolha dispensa a etapa de verificação da igualdade de variâncias (Fligner-Killeen), simplificando o fluxo de análise. Essa conveniência, no entanto, em caso de igualdade de variâncias, significa uma perda de poder de teste que pode ser especialmente relevante em amostras com poucas observações, como é o caso deste estudo. Por esse motivo, optou-se por verificar explicitamente a igualdade de variâncias e aplicar o teste mais adequado a cada caso, em vez de adotar Welch de forma indiscriminada.

\textbf{{[}PLACEHOLDER --- conciliar a formulação das hipóteses acima com o valor de referência
efetivamente utilizado nas chamadas do teste, e justificar a opção pela alternativa
bilateral.{]}}

As chamadas de \texttt{t.test()} foram executadas com o argumento \texttt{mu\ =\ delta\_star}, isto é, com
valor de referência igual a \(1\ \mathrm{kg/m^2}\); isso é refletido na linha
\texttt{alternative\ hypothesis:\ true\ difference\ in\ means\ is\ not\ equal\ to\ 1} das saídas obtidas.

\subsection{Descrição da Coleta e Consolidação dos Dados}\label{descriuxe7uxe3o-da-coleta-e-consolidauxe7uxe3o-dos-dados}

Por se tratar de um estudo retrospectivo, isto é, uma análise baseada em dados previamente coletados, a definição teórica do tamanho amostral e a estimação prévia do tamanho do efeito não puderam ser planejadas ativamente antes da coleta. Dessa forma, essas métricas serão reportadas e discutidas durante a apresentação dos resultados dos testes de hipótese. Esta subseção é dedicada exclusivamente ao procedimento de filtragem, tratamento e consolidação das bases de dados.

Os dados do semestre 2016-2 foram recuperados do arquivo \texttt{data/imc\_20162.csv}, enquanto os do semestre 2017-2 correspondem ao arquivo \texttt{data/CS01\_20172.csv}. Conforme especificado no enunciado, a base de 2016-2 inclui registros de turmas de graduação, e ambas as tabelas possuem divergências na estrutura das colunas. A consolidação dos dados exigiu as seguintes etapas de pré-processamento:

\begin{enumerate}
\def\labelenumi{\arabic{enumi}.}
\item
  \textbf{Filtragem dos dados de 2016-2:} isolamento dos alunos pertencentes à pós-graduação (\texttt{Course\ ==\ "PPGEE"}).
\item
  \textbf{Padronização dos nomes das colunas:} renomeação dos campos no arquivo de 2017-2 (\texttt{height.m} \(\rightarrow\) \texttt{Height.m} e \texttt{Sex} \(\rightarrow\) \texttt{Gender}) para alinhar com o esquema do arquivo de 2016-2.
\item
  \textbf{Conversão para tipo numerico}
\item
  \textbf{Cálculo do IMC:} criação da variável \texttt{BMI} em ambas as tabelas, aplicando a fórmula:
  \[\mathrm{BMI} = \frac{\mathrm{Peso}}{\mathrm{Altura}^2}\]
\item
  \textbf{Estruturação por subpopulações:} separação das observações por sexo e semestre para a criação dos conjuntos agregados que compõem os grupos de análise, alocados nas seguintes variáveis:

  \begin{itemize}
  \tightlist
  \item
    \texttt{fem\_2016}, \texttt{fem\_2017}: subpopulações femininas dos semestres 2016-2 e 2017-2, respectivamente.
  \item
    \texttt{masc\_2016}, \texttt{masc\_2017}: subpopulações masculinas dos semestres 2016-2 e 2017-2, respectivamente.
  \item
    \texttt{fem}, \texttt{masc}: agregados das subpopulações femininas e masculinas, englobando ambos os semestres.
  \item
    \texttt{todos}: tabela consolidada com a totalidade dos dados, classificados por sexo e semestre.
  \end{itemize}
\end{enumerate}

\section{Análise Exploratória dos Dados}\label{anuxe1lise-exploratuxf3ria-dos-dados}

A fim de compreender o comportamento dos dados, foi realizada uma exploração inicial por meio de um gráfico de \emph{boxplot} para comparar a distribuição do IMC entre as subpopulações de cada semestre.

\begin{figure}[H]

{\centering \includegraphics[width=0.7\linewidth]{ec1_report_files/figure-latex/grafico-explo-dados-1} 

}

\caption{Distribuição do IMC por Sexo e Semestre}\label{fig:grafico-explo-dados}
\end{figure}

Em seguida, os principais parâmetros estatísticos (tamanho amostral \(n\), média e desvio-padrão) de cada subcategoria foram calculados e resumidos na Tabela \ref{tab:tabela-resumo}.

\begin{table}[H]
\centering
\caption{\label{tab:tabela-resumo}Estatísticas descritivas do IMC por subamostra}
\centering
\begin{tabular}[t]{l|r|r|r}
\hline
Grupo & n & Média & Desvio Padrão\\
\hline
Feminino 2016-2 & 7 & 21.08 & 2.42\\
\hline
Feminino 2017-2 & 4 & 18.45 & 1.60\\
\hline
Masculino 2016-2 & 21 & 24.94 & 4.32\\
\hline
Masculino 2017-2 & 21 & 24.29 & 3.44\\
\hline
\end{tabular}
\end{table}

\section{Análise Estatística}\label{anuxe1lise-estatuxedstica}

\subsection{Validação das premissas do modelo}\label{validauxe7uxe3o-das-premissas-do-modelo}

Em cada caso, as premissas do teste foram verificadas na ordem
normalidade \(\rightarrow\) igualdade de variâncias \(\rightarrow\) independência. Os p-valores
do teste de Durbin-Watson são obtidos por reamostragem; por isso a semente aleatória foi
fixada em \texttt{set.seed(1234)}, garantindo a reprodutibilidade dos valores reportados.

\textbf{{[}PLACEHOLDER: discussão do Caso 1{]}}
\textbf{{[} verificar a necessidade do teste de independecia {]}}

\subsubsection{Normalidade}\label{normalidade}

O teste de Shapiro-Wilk foi realizado para cada uma das subamostras em análise. Os resultados seguem na Tabela \ref{tab:st-table}.

\begin{table}

\caption{\label{tab:st-table}Tabela 2 : Resultado do p-valor para o teste de Shapiro-Wilk }
\centering
\begin{tabular}[t]{l|l|r}
\hline
Genero & Semestre & p-valor\\
\hline
F & 2016-2 & 0.467\\
\hline
M & 2016-2 & 0.127\\
\hline
F & 2017-2 & 0.037\\
\hline
M & 2017-2 & 0.621\\
\hline
\end{tabular}
\end{table}

Os resultados indicam que as amostras masculinas (2016-2 e 2017-2) e a amostra feminina de 2016-2 apresentam comportamento compatível com a distribuição normal (\(p\text{-valor} > 0{,}05\)), não permitindo a rejeição da hipótese nula de normalidade.

Por outro lado, a amostra feminina do semestre 2017-2 (fem\_2017) apresentou um \(p\text{-valor}\) de \(0{,}036\), ficando abaixo do nível de significância adotado (\(\alpha = 0{,}05\)). Isso indica não há evidências estatísticas suficientes para assumir a normalidade desse subgrupo. No entanto, por se tratar de um \(p\text{-valor}\) marginalmente próximo ao limiar de \(0{,}05\), é plausível supor que a presença de algum ponto discrepante (outlier) ou a assimetria da amostra esteja afetando o teste. Ao analisar o boxplot da Figura \ref{fig:grafico-explo-dados}, não se observa um outlier extremo para este subgrupo, contudo nota-se que o valor máximo está relativamente afastado da média. Por esse motivo, fez-se necessária uma investigação complementar por meio do gráfico quantil-quantil (Q-Q plot).

\begin{figure}

{\centering \includegraphics[width=0.7\linewidth]{ec1_report_files/figure-latex/qq-fem2017-1} 

}

\caption{Gráficos quantil-quantil do IMC do sub grupo feminino 2017}\label{fig:qq-fem2017}
\end{figure}

O gráfico quantil-quantil (Q-Q plot) indica que a violação da premissa de normalidade para a amostra feminina 2017-2 é provocada por apenas um único valor discrepante. Dessa forma, é razoável assumir a premissa de normalidade para as análises subsequentes, uma vez que o teste \(t\) de Student é robusto a pequenas desviações da distribuição normal.

Tal assunção é interessante, uma vez que o teste \(t\) de Student tradicional apresenta a vantagem de proporcionar maior poder de teste em relação ao teste de Welch, quando a premissa de homogeneidade de variâncias é de fato satisfeita.

\subsubsection{Igualdade de Variância}\label{igualdade-de-variuxe2ncia}

\begin{table}[H]
\centering
\caption{\label{tab:tabela-var}Resultado do teste de Fligner-Killeen para igualdade de variâncias entre semestres}
\centering
\begin{tabular}[t]{l|r|r}
\hline
Grupo de Comparação & Estatística (Chi²) & p-valor\\
\hline
Feminino (2016-2 vs 2017-2) & 0.709 & 0.400\\
\hline
Masculino (2016-2 vs 2017-2) & 0.083 & 0.774\\
\hline
\end{tabular}
\end{table}

De acordo com os resultados do teste de Fligner-Killeen apresentados na Tabela \ref{tab:tabela-var}, tanto a subpopulação feminina (\(p\text{-valor} > 0{,}05\)) quanto a masculina (\(p\text{-valor} > 0{,}05\)) mantiveram a premissa de homocedasticidade. Com a confirmação da igualdade de variâncias e da normalidade, o teste \(t\) de Student para amostras independentes com variâncias iguais é o procedimento adequado para a comparação das médias do IMC entre os semestres 2016-2 e 2017-2.

\subsubsection{\texorpdfstring{Efeito minimamente relevante (\(\delta^*\))}{Efeito minimamente relevante (\textbackslash delta\^{}*)}}\label{efeito-minimamente-relevante-delta}

O valor do menor efeito de interesse prático foi fixado \emph{a priori}, antes do exame dos dados,
em \(\delta^* = 1\ \mathrm{kg/m^2}\), tal valor foi adotado nos 2 casos analisados,
por ser a relevância prática uma propriedade da grandeza medida e não da dispersão de cada
subamostra. Essa escolha não pretende representar um limiar de relevância clínica ou prática específico, mas sim estabelecer uma referência comum e consistente, que permite comparar a sensibilidade do teste entre os dois grupos em igualdade de condições, evitando que diferenças nos resultados da discussão de potência decorram da escolha de valores de referência distintos para cada subpopulação, e não de diferenças reais em seus tamanhos amostrais ou variabilidades.

\subsubsection{Nível de significância e tamanho amostral}\label{nuxedvel-de-significuxe2ncia-e-tamanho-amostral}

Os testes foram executados com \texttt{conf.level\ =\ 0.95}, o que corresponde a um nível de
significância \(\alpha = 0{,}05\).

\subsection{Teste de Hipóteses}\label{teste-de-hipuxf3teses}

Os testes de hipótese foram realizados seguindo as hipóteses definidas na seção de Objetivo do Experimento e Hipóteses a serem Testadas. Os testes foram realizados para obter informações sobre a inferência da diferença da média do IMC para as subpopulações femininas 2016-2 vs.~2017-2, e repetidos para as subpopulações masculinas 2016-2 vs.~2017-2. Portanto, esta seção se divide em dois casos: subpopulação masculina e subpopulação feminina

\subsubsection{Caso 1: Subpopulação Masculina}\label{caso-1-subpopulauxe7uxe3o-masculina}

\begin{verbatim}
## 
##  Two Sample t-test
## 
## data:  masc_2016$BMI and masc_2017$BMI
## t = 0.53979, df = 40, p-value = 0.5923
## alternative hypothesis: true difference in means is not equal to 0
## 95 percent confidence interval:
##  -1.784943  3.085836
## sample estimates:
## mean of x mean of y 
##  24.93595  24.28551
\end{verbatim}

O teste resultou em \(p = 0{,}5923\), sendo o intervalo de confiança de 95\% para a diferença entre as médias: \([-1{,}7849;\ 3{,}0858]\ \mathrm{kg/m^2}\), e médias amostrais de \(24{,}9360\ \mathrm{kg/m^2}\) (2016-2) e \(24{,}2855\ \mathrm{kg/m^2}\) (2017-2). Como o intervalo de confiança contém o valor zero e, o \(p\)-valor obtido é maior que \(\alpha = 0{,}05\), falhou-se em rejeitar a hipótese nula, não havendo evidências estatísticas suficientes para afirmar que o IMC médio da subpopulação masculina difere entre os semestres 2016-2 e 2017-2.

\subsubsection{Caso 2: População Feminina}\label{caso-2-populauxe7uxe3o-feminina}

\begin{verbatim}
## 
##  Two Sample t-test
## 
## data:  fem_2016$BMI and fem_2017$BMI
## t = 1.9308, df = 9, p-value = 0.08556
## alternative hypothesis: true difference in means is not equal to 0
## 95 percent confidence interval:
##  -0.4527037  5.7283762
## sample estimates:
## mean of x mean of y 
##  21.08443  18.44660
\end{verbatim}

O teste resultou em \(p = 0{,}08556\), sendo o intervalo de confiança de 95\% para a diferença entre as médias: \([-0{,}4527;\ 5{,}7284]\ \mathrm{kg/m^2}\), com médias amostrais de \(21{,}0844\ \mathrm{kg/m^2}\) (2016-2) e \(18{,}4466\ \mathrm{kg/m^2}\) (2017-2). Como o intervalo de confiança contém o valor zero, e o \(p\)-valor obtido, embora relativamente próximo de \(\alpha = 0{,}05\), ainda é maior que esse limiar, falhou-se em rejeitar a hipótese nula, não havendo evidências estatísticas suficientes para afirmar que o IMC médio da subpopulação feminina difere entre os semestres 2016-2 e 2017-2.

\subsection{Resultados e Discussões}\label{resultados-e-discussuxf5es}

\subsubsection{Discussão sobre o Teste de Hipóteses}\label{discussuxe3o-sobre-o-teste-de-hipuxf3teses}

O objetivo deste experimento era verificar se existia diferença significativa no IMC médio entre duas populações de estudantes: alunos de pós-graduação do PPGEE-UFMG nos semestres 2016-2 e 2017-2.

Os testes foram realizados sobre essas populações, separadas por gênero. Em ambos os casos, falhou-se em rejeitar a hipótese nula, não havendo evidências estatísticas suficientes para afirmar que o IMC médio difere entre os semestres, tanto para a subpopulação masculina quanto para a feminina. No entanto, devido ao tamanho reduzido de cada subamostra, é interessante realizar uma análise de poder do teste, que pode ajudar a avaliar a possibilidade de um Erro Tipo II, isto é, a hipótese de que se falhou em rejeitar \(H_0\) mesmo havendo, de fato, uma diferença real entre as populações.Por esse motivo, a seção a seguir traz essa discussão à tona.

\subsubsection{Discussão sobre o Poder do Teste}\label{discussuxe3o-sobre-o-poder-do-teste}

É convecioanl definir o poder estatístico em 80\% (\(1 - \beta = 0{,}80\)). Contudo, por se tratar de um estudo retrospectivo, em que os dados já haviam sido previamente coletados, não houve controle sobre o tamanho amostral para garantir essa potência \emph{a priori}.

Para poder levar a cabo o calculo de poder de teste devemos conhecer a variância populacional \(\sigma^2\), como a mesma é desconhecida, é necessário estimá-la a partir dos dados amostrais. Assumindo a premissa de igualdade de variância entre as subpopulações (\(\sigma_1^2 \approx \sigma_2^2\)), confirmada pelo teste de Fligner-Killeen, utiliza-se o estimador ponderado da variância agrupada (\(S_p^2\)):

\begin{equation}
S_p^2 = \frac{(n_1 - 1)S_1^2 + (n_2 - 1)S_2^2}{n_1 + n_2 - 2}
\label{eq:var-agrupada}
\end{equation}

em que \(S_1^2\) e \(S_2^2\) são as variâncias amostrais das subpopulações de 2016-2 e 2017-2, e \(n_1\), \(n_2\) os respectivos tamanhos amostrais. O desvio-padrão agrupado (\(S_p = \sqrt{S_p^2}\)) é então utilizado como parâmetro de dispersão no cálculo do poder do teste.

\begin{verbatim}
## Desvio Padrão - masculino - : 3.905
\end{verbatim}

\begin{verbatim}
## Desvio Padrão - femenino - : 2.18
\end{verbatim}

\begin{Shaded}
\begin{Highlighting}[]
\CommentTok{\# Poder do teste para detectar delta* = 1, dado o n e Sp já calculados}
\NormalTok{poder\_masc }\OtherTok{\textless{}{-}} \FunctionTok{power.t.test}\NormalTok{(}
  \AttributeTok{n =}\NormalTok{ n1\_masc,}
  \AttributeTok{delta =} \DecValTok{1}\NormalTok{,}
  \AttributeTok{sd =}\NormalTok{ sd\_pooled\_masc,}
  \AttributeTok{sig.level =} \FloatTok{0.05}\NormalTok{,}
  \AttributeTok{type =} \StringTok{"two.sample"}\NormalTok{,}
  \AttributeTok{alternative =} \StringTok{"two.sided"}
\NormalTok{)}\SpecialCharTok{$}\NormalTok{power}

\CommentTok{\#Usamos como caso balanceado por simplicidade}
\NormalTok{poder\_fem }\OtherTok{\textless{}{-}} \FunctionTok{power.t.test}\NormalTok{(}
  \AttributeTok{n =}\NormalTok{ n1\_fem,  }
  \AttributeTok{delta =} \DecValTok{1}\NormalTok{,}
  \AttributeTok{sd =}\NormalTok{ sd\_pooled\_fem,}
  \AttributeTok{sig.level =} \FloatTok{0.05}\NormalTok{,}
  \AttributeTok{type =} \StringTok{"two.sample"}\NormalTok{,}
  \AttributeTok{alternative =} \StringTok{"two.sided"}
\NormalTok{)}\SpecialCharTok{$}\NormalTok{power}

\FunctionTok{cat}\NormalTok{(}\StringTok{"Poder para detectar delta* = 1 (masculino):"}\NormalTok{, }\FunctionTok{round}\NormalTok{(poder\_masc, }\DecValTok{3}\NormalTok{), }\StringTok{"}\SpecialCharTok{\textbackslash{}n}\StringTok{"}\NormalTok{)}
\end{Highlighting}
\end{Shaded}

\begin{verbatim}
## Poder para detectar delta* = 1 (masculino): 0.125
\end{verbatim}

\begin{Shaded}
\begin{Highlighting}[]
\FunctionTok{cat}\NormalTok{(}\StringTok{"Poder para detectar delta* = 1 (feminino):"}\NormalTok{, }\FunctionTok{round}\NormalTok{(poder\_fem, }\DecValTok{3}\NormalTok{), }\StringTok{"}\SpecialCharTok{\textbackslash{}n}\StringTok{"}\NormalTok{)}
\end{Highlighting}
\end{Shaded}

\begin{verbatim}
## Poder para detectar delta* = 1 (feminino): 0.121
\end{verbatim}

Observa-se que, em ambos os casos, o arranjo amostral disponível (com \(n_{masc}=42\) e \(n_{fem}=11\)) não teria sensibilidade suficiente para detectar, com poder de 80\%, uma diferença da magnitude que se definiu como referência mínima de interesse (\(\delta^*=1\ \mathrm{kg/m^2}\)). De fato, ao calcular diretamente o poder do teste para detectar especificamente \(\delta^*=1\ \mathrm{kg/m^2}\), obtém-se um poder de apenas 12,5\% (masculino) e 12,1\% (feminino), valores bem abaixo do limiar convencional de 80\% considerado adequado.

\section{Conclusões e Recomendações}\label{conclusuxf5es-e-recomendauxe7uxf5es}

Os resultado extraidos a parti do poder do teste nos diz que falha em rejeitar \(H_0\) em ambas as subpopulações não deve ser interpretada como evidência de ausência de diferença real entre os semestres. Uma vez que a probabilidade de cometer um erro do tipo 2 é aproximadamente de 88\% para ambas subpopulações para o delta definido.

\subsection{Sugestões de melhoria do experimento}\label{sugestuxf5es-de-melhoria-do-experimento}

\subsubsection{Menor Diferença Detectável}\label{menor-diferenuxe7a-detectuxe1vel}

Considerando, o baixo poder de teste calculado na seçao anterior, aqui se pretende discutir:

\begin{enumerate}
\def\labelenumi{\arabic{enumi}.}
\tightlist
\item
  o numero de observaçoes necessarios para atingir o poder de teste de 80\%;
\item
  a menor diferença detectavel para um poder de teste de 80\% e o nível de significância em \(\alpha = 0{,}05\), considerando o arrajo amostral disponivel.
\end{enumerate}

A função empregada para esse cálculo, porém, depende do balanceamento amostral de cada subpopulação: para a subpopulação masculina, balanceada (\(n_{2016} = n_{2017} = 21\)), utiliza-se \texttt{power.t.test()}; para a feminina, desbalanceada (\(n_{2016} \neq n_{2017}\)), utiliza-se \texttt{pwr.t2n.test()} do pacote \texttt{pwr}, que admite tamanhos amostrais distintos entre os dois grupos.

\textbf{Subpopulação Masculina}

\begin{verbatim}
## Número de observações necessário (delta* = 1) - feminino: 240.3
\end{verbatim}

\begin{verbatim}
## Menor diferença detectável - masculina: 3.46 kg/m²
\end{verbatim}

\textbf{Subpopulação Feminina}

\begin{verbatim}
## Número de observações necessário (delta* = 1) - feminino: 75.6
\end{verbatim}

\begin{verbatim}
## Menor diferença detectável - feminina: 4.3 kg/m²
\end{verbatim}

Essa limitação reforça a recomendação de que estudos futuros sobre o tema considerem um planejamento amostral prévio, com cálculo de tamanho de amostra baseado em um \(\delta^*\) definido antes da coleta, de modo a garantir poder adequado para detectar efeitos de magnitude prática relevante.

\subsection{Declaração sobre o Uso de IA Generativa e Tecnologias Assistidas por IA no Processo de Elaboração do Relatório}\label{declarauxe7uxe3o-sobre-o-uso-de-ia-generativa-e-tecnologias-assistidas-por-ia-no-processo-de-elaborauxe7uxe3o-do-relatuxf3rio}

Durante a preparação deste trabalho, os autores utilizaram o Claude (Anthropic) para redigir a versão inicial das seções descritivas deste relatório, a partir dos resultados registrados como comentários no script de análise \texttt{processamento\_separado.R}. Após utilizar essa ferramenta, os autores revisaram e editaram o conteúdo conforme necessário e assumem total responsabilidade pelo conteúdo da publicação.

Também foi utilizada IA para auxiliar na formatação do texto, incluindo a inclusão de bibliotecas que ajudam na visualização de gráficos e tabelas apresentados no relatório. Da mesma forma, foi solicitado apoio para revisar o \emph{chunk} de \emph{setup} do R Markdown, de modo a garantir a instalação das bibliotecas necessárias para que o arquivo compilasse sem erros.

Por fim, na seção sobre Menor Diferença Detectável, foi questionado qual função poderia ser utilizada para realizar o teste de poder no caso de populações não balanceadas, uma vez que esse caso não foi discutido (em R) durante as aulas.

\end{document}
